\section{Discussion}\label{sec:discussion}

\subsection{What the divergences teach}
The reproducibility study is, to our knowledge, the first time a constructed language's founding text has been tested against its own stated procedure. The result---the Manual is an algorithm for what things mean and who owns them, and a hand for how they sound---is informative in both directions. It shows that determinism is achievable for the semantically loaded decisions: no reader of the Manual could complain that \emph{kala} or \emph{duž} or \emph{peleek} were chosen by taste. And it locates the residue of taste precisely, in the filling-in of shared words and the counting of ballots. A second edition of the Manual could close the gap in one of two ways: by adopting the engine's fully specified fill-in rule and accepting \emph{silk} and \emph{penki}, or by stating explicitly that the Domain Rule takes precedence for the sea vocabulary and that Hanse words keep the vowel length of the family that ranks first in the queue at the time of derivation. Either would make Baltmeri deterministic to the last letter; our recommendation is the first, because a language that belongs to the sea should not privilege how the sea sounds in Estonian.

\subsection{Equity, measured}
The family-balance result (Table~\ref{tab:balance}) is the design's central empirical claim. Germanic, with three members and roughly 110 million first-language speakers, supplies the fewest forms; Baltic, with two members and under five million speakers, supplies the most. This is not a bias toward the small: it is what happens when a shared word is counted once per family and when the tie-break geography is the sea's rather than the speaker count's. Interslavic's equal weighting of six sub-branches \citep{merunka2019} produced the same effect within one family; Baltmeri shows it can be produced across four. The obverse is that Latvian, by its early place in the queue, decided fourteen ties before the Rota Rule pushed it to the back; a language's influence in Baltmeri is a function of where its coast lies, which is the design's intention and its most contestable feature.

\subsection{The Baltic as a subject}
We began from the observation that the sea has no name of its own. Baltmeri gives it one, and a way of talking about itself in which \emph{Meri er men} ``the sea is ours'' means what a citizen of the Baltic would want it to mean: not possession but belonging. The legal and institutional turn toward recognising ecosystems as persons \citep{stone1972,teawatupua2017,spain2022,kauffman2021} requires guardians who speak for the entity; \citet{thiel2024} has proposed exactly such a body for this sea. A guardian who speaks \emph{for} a sea still speaks in a national language. A guardian who can also speak \emph{as} the sea---in a language generated from all its shores and owned by none---has something new to say, and the everyday users of a translator that shows its work are, in a small way, practising the bioregional citizenship Thiel describes. We do not claim that a constructed language will restore oxygen to the Gotland Deep. We claim that identity precedes care, that \citet{gotz2016} is right that Baltic identity has not followed Baltic institutions, and that a language is one of the few instruments---\citet{anderson2006} showed the print-language was the decisive one---by which a community can imagine itself into being. The environmental humanities have given us the words for a sea that is not outside us \citep{neimanis2017,steinberg2015,haraway2016}; Baltmeri is an attempt to give the sea words of its own.

\subsection{Limitations}
The nine ``first dictionary translations'' on which every derivation depends are a human or model judgement, and lexical choice upstream can change a word downstream: had we given Russian \emph{ходить под парусом} rather than \emph{плыть} for ``sail'', the ballot would differ. The Manual's queue is a convention that does not exactly track geography (\S\ref{sec:sources}); the choice of Saint Petersburg rather than Kaliningrad as Russia's coastal city is a political decision the text does not acknowledge. The seed lexicon is small and skewed toward the sea; a Leipzig--Jakarta list of core meanings \citep{tadmor2010} would be the natural next seed, and would test the procedure where families disagree most. The translator's analysis stage is probabilistic and can misanalyse a sentence even though every word it then receives is derived deterministically; the legend is faithful to the trace, not to the analysis. Baltmeri has no speakers, and intelligibility to speakers of the nine languages---the property Interslavic has measured at 84\%---has not been tested. Finally, the language is a proposition, not a claim on anyone: the nine languages of the shore are not diminished by a tenth that is made of them.

\subsection{Future work}
Three extensions follow directly. A crowd-sourced dictionary in which speakers of the nine languages supply and dispute source forms would make the upstream judgement communal rather than editorial, and would let the ``first derivation wins'' convention become a public record of who coined what. Cloze-based intelligibility studies on the Interslavic model would test whether Baltmeri is, as designed, passively readable around the sea. And a spoken form---the translator already voices Baltmeri with a multilingual model left to decide how it sounds---invites a phonetic study of what accent nine coastlines produce when no one is in charge.

\section{Conclusion}
Baltmeri is a language with a procedure where other languages have a history. Its words are decided by the nine languages of the Baltic shore voting as four equal families, with ties broken by the geography of the sea itself; its grammar is the set of endings those families could agree on; its lexicon can be extended by anyone with nine dictionaries and can be audited by anyone with the Manual. Implemented as a deterministic engine, the procedure reproduces the meaning-bearing decisions of its founding text and exposes, in the remaining half, exactly where a human hand intervened. Made usable through a translator in which a language model analyses and explains but never decides, it lets the citizens of the Baltic hear their sea say, in words that belong to all of them, \emph{meri er men}.

\section*{Data and code availability}
The Manual, the engine, the seed lexicon with all source forms, the 117 tests, and the translator are open source at \url{https://github.com/martinkallstrom/baltmeri-translator}. The translator is deployed at \url{https://baltmeri-translator.kindship-ai.workers.dev}. Every figure in \S\ref{sec:eval} is produced by scripts in the repository from the engine's own output.

\section*{Acknowledgements}
The authors thank the seals of Gotland for their patience.
